\section{Guía del clasificador: Classifier Guidance}

\subsection{De la distribución conjunta a la distribución condicional}

El objetivo de la generación condicional es aprender la distribución condicional $p(x|y)$, donde $y$ es la información condicional (por ejemplo, una etiqueta de clase). El insight clave radica en la relación simple que existe entre la función score de la distribución conjunta y la de la distribución condicional:

$$
\nabla_{x_t}\log p(x_t, y) = \nabla_{x_t}\log p(x_t) + \nabla_{x_t}\log p(y|x_t)
$$

Esto se debe a que $p(x_t, y) = p(x_t)\cdot p(y|x_t)$; al tomar el logaritmo y calcular el gradiente respecto a $x_t$, el término $\log p(y)$ desaparece porque es constante en relación con $x_t$.

\begin{importantbox}{Fórmula central de Classifier Guidance}
En cada paso de eliminación de ruido de los modelos de difusión, se modifica la función score:
$$
\tilde{\nabla}_{x_t}\log p(x_t, y) = \underbrace{\nabla_{x_t}\log p(x_t)}_{\text{score incondicional (modelo preentrenado)}} + \underbrace{\nabla_{x_t}\log p(y|x_t)}_{\text{guía por gradiente del clasificador}}
$$
Correspondientemente, al mapear esto a la predicción de ruido, la regla de actualización es:
$$
\tilde{\epsilon}(x_t, t) = \epsilon_\theta(x_t, t) - \sqrt{1-\bar\alpha_t}\;\nabla_{x_t}\log p_\phi(y|x_t)
$$
donde $p_\phi(y|x_t)$ es un clasificador entrenado sobre datos con ruido.
\end{importantbox}

\subsection{Entrenamiento del clasificador de ruido}

Aquí hay un detalle crucial: el clasificador $p_\phi(y|x_t)$ no es un clasificador de imágenes convencional, sino que debe ser capaz de manejar puntos intermedios $x_t$ con \textbf{niveles arbitrarios de ruido}. El método de entrenamiento es el siguiente:

\begin{enumerate}
    \item Para cada par $(x_0, y)$ en el conjunto de datos, muestrear aleatoriamente un paso de tiempo $t$.
    \item Obtener la imagen con ruido $x_t$ mediante el proceso hacia adelante $q(x_t|x_0)$.
    \item Entrenar al clasificador para predecir la probabilidad de la clase $y$ utilizando $x_t$ y $t$ como entradas.
\end{enumerate}

\begin{warningbox}{El clasificador debe manejar entradas con ruido}
¡No se puede utilizar directamente un clasificador ImageNet entrenado en imágenes limpias! El clasificador debe entrenarse sobre $x_t$ con diversos niveles de ruido; de lo contrario, la señal de gradiente será completamente confiable en los pasos de alto ruido. Esto representa un costo de entrenamiento adicional para Classifier Guidance.
\end{warningbox}

\subsection{Control de la intensidad de la guía}

Mediante la introducción de un peso $w$ se controla la intensidad de la condición:

$$
\tilde{\epsilon}(x_t, t) = \epsilon_\theta(x_t, t) - w\cdot\sqrt{1-\bar\alpha_t}\;\nabla_{x_t}\log p_\phi(y|x_t)
$$

\begin{itemize}
    \item $w = 0$: Generación incondicional pura.
    \item $w = 1$: Generación condicional estándar (correspondiente al score de la distribución conjunta).
    \item $w > 1$: Efecto condicional aumentado —la salida es más fiel a la condición, pero disminuye la diversidad.
\end{itemize}

\begin{knowledgebox}{Comprensión intuitiva de la intensidad de la guía}
Desde la perspectiva de las trayectorias, la trayectoria de eliminación de ruido del modelo de difusión incondicional va de $x_T$ a $x_0$ siguiendo el campo de score incondicional. Agregar el término de gradiente del clasificador equivale a \textbf{desviar} la trayectoria en cada paso hacia la región del manifold de datos que satisface la condición $y$. Aumentar $w$ hace que la desviación sea más intensa, produciendo resultados más "típicos" pero con menor diversidad.
\end{knowledgebox}

\subsection{Ventajas y limitaciones de Classifier Guidance}

\begin{table}[H]
\centering
\begin{tabular}{p{6cm}p{6cm}}
\toprule
\textbf{Ventajas} & \textbf{Limitaciones} \\
\midrule
Reutilización de modelos incondicionales preentrenados & Requiere entrenamiento adicional de un clasificador de ruido \\
Control flexible de la intensidad de la guía & Se debe ejecutar dos redes durante la inferencia \\
No requiere reentrenamiento del modelo de difusión & Los gradientes del clasificador pueden ser inestables \\
Derivación matemática clara & Limitado a tipos de condición para los que está disponible un clasificador \\
\bottomrule
\end{tabular}
\caption{Comparación de ventajas y desventajas de Classifier Guidance}
\end{table}

Estas limitaciones impulsaron directamente la propuesta de Classifier-Free Guidance.

\subsection{Resumen de la sección}

La idea central de Classifier Guidance es: aprovechar la descomposición bayesiana para descomponer la generación condicional en un score incondicional más el gradiente del clasificador. Este insight revela una ventaja única de los modelos de difusión frente a las GAN y los VAE: pueden lograr orientación condicional mediante la manipulación del score function \textbf{sin reentrenar} el modelo generador. Esta paradigma de manipulación de scores se mantendrá presente en todo el contenido posterior del curso.

\newpage
